On the complete GB1 four-site landscape, with five draws per cell, we found no detectable benefit from pairwise epistasis features under random training (Spearman \rhval{main/one-hot-+-pairwise-ridge/rand_2000/spearman/mean:3} versus \rhval{main/one-hot-ridge-reimplemented/rand_2000/spearman/mean:3} for the additive model at $N{=}2000$) and a large loss when training on single and double mutants and testing on higher-order variants (\rhval{main/one-hot-+-pairwise-ridge/dbl_2000/spearman/mean:3} versus \rhval{main/one-hot-ridge-reimplemented/dbl_2000/spearman/mean:3}). A Hamming-kernel GP matched additive ridge at small N and fell behind at $N{=}2000$. A small CNN had the best mean Spearman at $N\ge96$, but a shared log-transformed target brought the linear models and the GP to within a narrow band of it at $N{=}384$, which is consistent with part of the CNN's advantage on the raw target being a learned output nonlinearity (the gap narrowed but did not vanish at $N{=}2000$). Of the five hypotheses, H5 is supported, H2 is supported in the mean but statistically inconclusive, and H1, H3 and H4 are refuted as stated (H4 only at $N\le96$). These are small-scale results on one landscape with simple reimplemented baselines; they should be read as a calibrated reference point, not as a ranking of methods.
